% Part: incompleteness
% Chapter: incompleteness-provability
% Section: rosser-thm

\documentclass[../../../include/open-logic-section]{subfiles}

\begin{document}

\olfileid{inc}{inp}{ros}

\olsection{Teorema Rosser}

Apa bukti G\"odel bisa diowahi supaya ngasilake asil kang luwih kuwat,
kanthi ngganti ``$\omega$-konsisten'' dadi mung ``konsisten''?
Wangsulane ``bisa,'' nganggo trik kang ditemokake Rosser. Trik
Rosser yaiku nggunakake predikat !!{derivability} kang
``diowahi'' $\ORProv_T(y)$ tinimbang $\OProv[\Th{T}](y)$.

\begin{thm}
\ollabel{thm:rosser}
Upamana $\Th{T}$ teori konsisten sembarang kang !!{axiomatizable}
lan ngembangake $\Th{Q}$. Mula $\Th{T}$ ora jangkep.
\end{thm}

\begin{proof}
Elinga yen $\OProv[\Th{T}](y)$ ditemtokake minangka
$\lexists[x][\OPrf[\Th{T}](x, y)]$, ing ngendi
$\OPrf[\Th{T}](x, y)$ makili relasi kang bisa diputusake
lan bener yen lan mung yen $x$ iku wilangan G\"odel saka
!!a{derivation} kanggo !!{sentence} kang nduweni wilangan
G\"odel~$y$. Relasi antarane $x$ lan~$y$ kang bener yen $x$~iku
wilangan G\"odel saka sawijining \emph{bantahan} tumrap ukara kang
nduweni wilangan G\"odel~$y$ uga bisa diputusake. Tetepna
$\fn{not}(x)$ minangka fungsi rekursif primitif mangkene: yen $x$
kode saka formula $!A$, $\fn{not}(x)$ iku kode saka $\lnot !A$.
Mula $\Refut[\Th{T}](x, y)$ bener yen lan mung yen
$\Prf[\Th{T}](x, \fn{not}(y))$. Tetepna
$\ORefut[\Th{T}](x, y)$ kanggo makili relasi iki. Mula,
yen $\Th{T} \Proves \lnot !A$ lan $\delta$ iku !!{derivation}
kang cocog, $\Th{Q} \Proves \ORefut[\Th{T}](\gn{\delta}, \gn{!A})$.
Kita nemtokake $\ORProv[\Th{T}](y)$ minangka
\[
\lexists[x][(\OPrf[\Th{T}](x,y) \land \lforall[z][(z < x \lif \lnot
  \ORefut[\Th{T}](z,y))])].
\]
Kanthi kasar, $\ORProv[\Th{T}](y)$ nyatakake ``ana bukti kanggo
$y$ ing $\Th{T}$, lan ora ana bantahan kanggo~$y$ kang kodene luwih
cilik tinimbang kode bukti mau.'' Perbandhingan ing rumus iki adhedhasar ukuran angka kode bukti lan bantahan.
Yen $\Th{T}$ konsisten, $\ORProv[\Th{T}](y)$ bener tumrap wilangan kang
padha karo $\OProv[\Th{T}](y)$; nanging saka sudut pandang
\emph{provabilitas} ing~$\Th{T}$ (lan saiki kita ngerti yen bebener lan
provabilitas iku beda!) loro-lorone nduweni sipat kang beda. Yen
$\Th{T}$ \emph{ora} konsisten, loro-lorone \emph{ora} bener tumrap
wilangan kang padha!{} ($\ORProv[\Th{T}](y)$ kerep diwaca minangka
``$y$ bisa dibuktekake kanthi cara Rosser.'' Nanging, kaya kang
lagi dirembug, provabilitas Rosser dudu jinis provabilitas khusus---ing
teori kang ora konsisten, ana !!{sentence}s kang bisa dibuktekake
nanging ora bisa dibuktekake kanthi cara Rosser---mula sebutan iki
bisa mbingungake. Supaya ora bingung, bisa wae diwaca minangka
``$y$ iku shmovable,'' nganggo tembung guyon saka sumber)

Miturut lema titik tetep, ana ukara $!R_\Th{T}$ saengga
\begin{equation}
  \Th{Q} \Proves !R_\Th{T} \liff \lnot \ORProv[\Th{T}](\gn{!R_\Th{T}}).
  \ollabel{RT}
\end{equation}
Beda karo bukti \olref[1in]{thm:first-incompleteness}, ing kene kita
nyatakake yen $\Th{T}$ konsisten, $\Th{T}$ ora !!{derive}
$!R_\Th{T}$, lan $\Th{T}$ uga ora !!{derive} $\lnot !R_\Th{T}$.
(Kanthi tembung liya, kita ora mbutuhake panganggep
$\omega$-konsisten.)

Sepisan, ayo dituduhake yen $\Th{T} \Proves/ !R_{\Th{T}}$.
Upamana teori iki mbuktekake, mula ana !!a{derivation}
kanggo~$!R_{\Th{T}}$ saka~$T$; upamana $n$ wilangan G\"odele. Mula
$\Th{Q} \Proves \OPrf[\Th{T}](\num{n}, \gn{!R_{\Th{T}}})$, amarga
$\OPrf[\Th{T}]$ makili $\Prf[\Th{T}]$ ing~$\Th{Q}$.
Kajaba iku, kanggo saben $k < n$, $k$ dudu wilangan G\"odel saka
!!a{derivation} kanggo $\lnot !R_{\Th{T}}$, amarga $\Th{T}$
konsisten. Dadi kanggo saben $k < n$, $\Th{Q} \Proves \lnot
\ORefut[\Th{T}](\num{k}, \gn{!R_{\Th{T}}})$. Miturut
\olref[req][min]{lem:less-nsucc},
$\Th{Q} \Proves \lforall[z][(z < \num{n} \lif \lnot \ORefut[\Th{T}](z,
  \gn{!R_{\Th{T}}}))]$. Mula,
\[
\Th{Q} \Proves \lexists[x][(\OPrf[\Th{T}](x,\gn{!R_{\Th{T}}}) \land \lforall[z][(z
    < x \lif \lnot \ORefut[\Th{T}](z,\gn{!R_{\Th{T}}}))])],
\]
nanging iki persis $\ORProv[\Th{T}](\gn{!R_{\Th{T}}})$. Miturut
\olref{RT}, $\Th{Q} \Proves \lnot !R_{\Th{T}}$. Amarga $\Th{T}$
ngembangake $\Th{Q}$, uga $\Th{T} \Proves \lnot !R_{\Th{T}}$.
Kita wis nganggep yen $\Th{T} \Proves !R_{\Th{T}}$, mula $\Th{T}$
bakal ora konsisten, nalisir panganggep ing teorema.

Saiki ayo dituduhake yen $\Th{T} \Proves/ \lnot !R_{\Th{T}}$.
Upamana maneh teori iki mbuktekake, lan upamana $n$ wilangan G\"odel
saka !!a{derivation} kanggo~$\lnot !R_{\Th{T}}$. Mula
$\Refut[\Th{T}](n, \Gn{!R_{\Th{T}}})$ bener, lan amarga
$\ORefut[\Th{T}]$ makili $\Refut[\Th{T}]$ ing
$\Th{Q}$, $\Th{Q} \Proves
\ORefut[\Th{T}](\num{n}, \gn{!R_{\Th{T}}})$. Kita bakal nuduhake
yen $\Th{T}$ banjur ora konsisten amarga teori iki uga
!!{derive}~$!R_{\Th{T}}$. Amarga
\begin{align*}
\Th{Q} & \Proves !R_{\Th{T}} \liff \lnot \ORProv[\Th{T}](\gn{!R_{\Th{T}}}), 
\intertext{lan amarga $\Th{T}$ ngembangake~$\Th{Q}$, cukup dituduhake yen}
\Th{Q} & \Proves \lnot \ORProv[\Th{T}](\gn{!R_{\Th{T}}}).
\end{align*}
!!{sentence} $\lnot
\ORProv[\Th{T}](\gn{!R_{\Th{T}}})$, yaiku,
\begin{align*}
  \lnot & \lexists[x][(\OPrf[\Th{T}](x,\gn{!R_{\Th{T}}}) \land
    \lforall[z][(z < x \lif \lnot \ORefut[\Th{T}](z,\gn{!R_{\Th{T}}}))])],
  \intertext{ekuivalen sacara logis karo}
  & \lforall[x][(\OPrf[\Th{T}](x,\gn{!R_{\Th{T}}}) \lif
    \lexists[z][(z < x \land \ORefut[\Th{T}](z,\gn{!R_{\Th{T}}}))])].
\end{align*}
Kita ngandharake kanthi informal nganggo logika lan fakta-fakta kang
$\Th{Q}$ !!{derive}s. Upamana $x$ sembarang lan
$\OPrf[\Th{T}](x,\gn{!R_{\Th{T}}})$. Kita wis ngerti yen
$\Th{T} \Proves/ !R_{\Th{T}}$, mula kanggo saben $k$,
$\Th{Q} \Proves \lnot \OPrf[\Th{T}](\num{k}, \gn{!R_{\Th{T}}})$.
Mula, kanggo saben instansi $k$ kang dibutuhake, kita entuk $\eq/[x][\num{k}]$. Mligine,
mung instansi nol nganti n kang dibutuhake ing bukti iki, mula ora ana panganggone aturan omega. (a) kita nduweni $\eq/[x][\num{n}]$. Kita uga nduweni
$\lnot(\eq[x][\num{0}] \lor \eq[x][\num{1}] \lor \dots \lor
\eq[x][\num{n-1}])$ lan mula miturut
\olref[req][min]{lem:less-nsucc}, (b) $\lnot(x < \num{n})$.
Yen n nol, ora ana kasus kurang saka n; lema kurang-saka-nol menehi (b) langsung. Yen n positif, disjungsi kasus ing ndhuwur winates. Miturut \olref[req][min]{lem:trichotomy}, $\num{n} < x$.
Amarga $\Th{Q} \Proves
\ORefut[\Th{T}](\num{n}, \gn{!R_{\Th{T}}})$, kita nduweni
$\num{n} < x \land
\ORefut[\Th{T}](\num{n}, \gn{!R_{\Th{T}}})$, lan saka iku
$\lexists[z][(z < x \land \ORefut[\Th{T}](z,\gn{!R_{\Th{T}}}))]$.
Amarga $x$ sembarang, kita entuk kaya kang dibutuhake:
\[
\lforall[x][(\OPrf[\Th{T}](x,\gn{!R_{\Th{T}}}) \lif
  \lexists[z][(z < x \land \ORefut[\Th{T}](z,\gn{!R_{\Th{T}}}))])].
\]
\end{proof}

\begin{prob}
Rong himpunan $A$ lan $B$ saka wilangan natural diarani
\emph{ora bisa dipisahake kanthi komputabel} yen ora ana himpunan
!!{decidable}~$X$ saengga $A \subseteq X$ lan
$B \subseteq \Complement{X}$ ($\Complement{X}$ iku komplemen,
$\Nat \setminus X$, saka~$X$). Upamana $\Th{T}$ ekstensi konsisten
kang !!{axiomatizable} saka $\Th{Q}$. Upamana $A$ himpunan wilangan
G\"odel saka !!{sentence}s kang bisa dibuktekake ing~$\Th{T}$ lan
$B$ himpunan wilangan G\"odel saka ukara kang bisa dibantah
ing~$\Th{T}$. Buktekna yen $A$ lan~$B$ ora bisa dipisahake
kanthi komputabel.
\end{prob}

\end{document}
